\documentclass[12pt]{letter}
\usepackage[margin=1.2in]{geometry}
\usepackage{hyperref}

\address{Independent Researcher\\Nanjing, China}
\signature{Huang Zhongchang}

\begin{document}

\begin{letter}{Editor\\Physical Review A}

\opening{Dear Editor,}

We submit ``Capacity-Normalized Bound on Non-Markovian Backflow from Finite Information Capacity'' for consideration as an Article in Physical Review A.

\textbf{What is new.}
The result provides an operationally accessible bound on non-Markovian backflow requiring only projective population measurements---no spectral densities, coupling strengths, or joint-state tomography. The bound follows from the pigeonhole principle applied to qubit toggle events under unidirectional Lindblad dynamics: each qubit can toggle $|0\rangle \to |1\rangle$ at most once, yielding $\mathcal{R} \equiv N_{\rm back}/C_F \leq q_S/q_E$ (for equal system and environment sizes). The Megier--Smirne--Vacchini entropic bound [Phys. Rev. Lett. 127, 030401 (2021)] applies to any CPTP dynamics but requires full tomography of the joint system-environment state; the present bound is narrower in domain but immediately measurable on current superconducting hardware.

\textbf{Why PRA.}
This result belongs to the foundations of open quantum systems, for which Physical Review A has long been the central venue. The Breuer--Laine--Piilo trace-distance measure and the Rivas--Huelga--Plenio CP-divisibility criterion---the two pillars of modern non-Markovianity theory---both appeared in Physical Review Letters, but the broader landscape of non-Markovian measures, operational characterizations, and collision-model analyses has been developed primarily in Physical Review A. The present work contributes directly to this tradition: a combinatorially derived bound that connects the finite information capacity of the physical substrate to an operationally accessible constraint on backflow.

\textbf{Relation to prior work.}
The Breuer--Laine--Piilo measure detects non-Markovianity via trace-distance revivals; Rivas, Huelga, and Plenio [PRL 105, 050403 (2010)] characterized it via CP-divisibility; Megier, Smirne, and Vacchini [PRL 127, 030401 (2021)] bound backflow entropically. Buscemi et al.\ [PRX Quantum 6, 020316 (2025)] classify revivals as causal or noncausal within the process-matrix framework. Our contribution is a combinatorially-derived, operationally simple bound that is complementary to MSV---narrower in domain (unidirectional Lindblad channels) but requiring only projective measurements.

Suggested referees: Bassano Vacchini (non-Markovian dynamics, co-author of the MSV bound), \'A.\ Rivas (CP-divisibility framework), Dariusz Chru\'sci\'nski (open quantum systems), and Fabio Costa (quantum causality).

\closing{Sincerely,}

\end{letter}
\end{document}
